% !TEX root = ../main.tex
%% Chapitre 17 — Probabilités sur un univers fini
%% Programme officiel MPSI/MP2I (BO spécial n°1 du 11/02/2021), deuxième semestre,
%% section « Probabilités » : A) univers, événements, variables aléatoires ;
%% espaces probabilisés finis ; conditionnelles (composées, totales, Bayes) ;
%% lois (uniforme, Bernoulli, binomiale), couples, marginales ; indépendance
%% d'événements et de variables (lemme des coalitions) ; B) espérance (transfert),
%% variance, écart type, covariance ; C) inégalités de Markov et de
%% Bienaymé-Tchebychev, loi faible des grands nombres.
%% Résultats vérifiés par verifications/chapitre_17.py.

\begin{ouverture}[binomiale]{17}{Probabilités sur un univers fini}{Partie 2 · Deuxième semestre}
  {Espaces probabilisés\enspace·\enspace Conditionnement\enspace·\enspace Lois\enspace·\enspace
   Indépendance\enspace·\enspace Espérance et variance\enspace·\enspace Concentration}
  \objectif{Modéliser une expérience aléatoire : univers fini, événements, variables aléatoires.}
  \objectif{Conditionner : probabilités composées, formule des probabilités totales, Bayes.}
  \objectif{Manier les lois usuelles : uniforme, Bernoulli, binomiale — et les couples de variables.}
  \objectif{Reconnaître et exploiter l'indépendance, d'événements comme de variables.}
  \objectif{Calculer espérances, variances et covariances (transfert, König-Huygens).}
  \objectif{Contrôler les écarts : Markov, Bienaymé-Tchebychev, loi faible des grands nombres.}
  \prerequis{Dénombrement : principe multiplicatif, combinaisons, binôme (chapitre 16).}
  \prerequis{Ensembles et applications (chapitre 1) ; sommes finies (chapitre 2).}
  \prerequis{Analyse asymptotique pour les comportements limites (chapitre 11).}
\end{ouverture}

%% =====================================================================
%%  L'ESSENTIEL
%% =====================================================================
\essentiel

\essSection{1}{Univers, événements, variables aléatoires}

\begin{definition}[Le vocabulaire probabiliste]
Une expérience aléatoire est modélisée par un \textbf{univers} fini
$\Omega$ (les issues possibles). Un \textbf{événement} est une partie de
$\Omega$, un événement \emph{élémentaire} un singleton. Le dictionnaire
ensembliste se traduit :
\begin{center}
\begin{tabularx}{0.9\linewidth}{@{}>{\centering\arraybackslash}X>{\centering\arraybackslash}X@{}}
\textbf{ensembles} & \textbf{événements}\\[0.4mm]
$A\cup B$ & $A$ ou $B$\\
$A\cap B$ & $A$ et $B$\\
$\overline A$ & événement contraire\\
$A\cap B=\varnothing$ & incompatibles (disjoints)\\
partition de $\Omega$ & système complet d'événements\\
\end{tabularx}
\end{center}
Une \textbf{variable aléatoire} est une application $X:\Omega\to E$ ; on
note $(X\in A)=\ens{\omega\mid X(\omega)\in A}$ et $(X=x)$ : les
événements s'écrivent désormais en termes de variables.
\end{definition}

\essSection{2}{Espaces probabilisés finis}

\begin{definition}[Probabilité]
Une \textbf{probabilité} sur $\Omega$ est une application
$P:\parties{\Omega}\to\intFF{0}{1}$ telle que $P(\Omega)=1$ et
$P(A\cup B)=P(A)+P(B)$ pour $A,B$ incompatibles. Elle est entièrement
déterminée par la \emph{distribution} $\bigl(P(\ens{\omega})\bigr)_{\omega\in\Omega}$ —
une famille de réels positifs de somme $1$ — via
$P(A)=\sum_{\omega\in A}P(\ens{\omega})$. La probabilité
\textbf{uniforme} donne le même poids à chaque issue :
$P(A)=\dfrac{|A|}{|\Omega|}$ — le dénombrement du chapitre 16 devient du
calcul de probabilités.
\end{definition}

\begin{propriete}[Règles de calcul]
$P(\overline A)=1-P(A)$ ;
$P(A\cup B)=P(A)+P(B)-P(A\cap B)$ ;
$P(A\setminus B)=P(A)-P(A\cap B)$ ; et la croissance :
$A\subseteq B\implies P(A)\leqslant P(B)$. (La formule du crible générale
est hors programme.)
\end{propriete}

\essSection{3}{Probabilités conditionnelles}

\noindent
\begin{minipage}[c]{0.55\linewidth}
\begin{definition}[Conditionnement]
Si $P(B)>0$, la \textbf{probabilité conditionnelle} de $A$ sachant $B$
est
\[
P(A\mid B)=P_B(A)=\frac{P(A\cap B)}{P(B)} ,
\]
et $P_B$ est elle-même une probabilité sur $\Omega$. Par convention,
$P(A\mid B)P(B)=0$ si $P(B)=0$.
\end{definition}
\begin{theoreme}[Les trois formules]
\emph{Composées} :
$P(A_1\cap\dots\cap A_n)
=P(A_1)\,P(A_2\mid A_1)\cdots P(A_n\mid A_1\cap\dots\cap A_{n-1})$.
\emph{Probabilités totales} : si $(B_i)$ est un système complet
d'événements,
$P(A)=\sum_i P(A\mid B_i)\,P(B_i)$.
\emph{Bayes} : si de plus $P(A)>0$,
$P(B_j\mid A)=\dfrac{P(A\mid B_j)P(B_j)}{\sum_iP(A\mid B_i)P(B_i)}$ —
remonter des effets aux causes.
\end{theoreme}
\end{minipage}\hfill
\begin{minipage}[c]{0.41\linewidth}
\centering
\begin{tikzpicture}[scale=0.92,line cap=round]
  \fill[otInk] (0,1.5) circle (1.6pt);
  \fill[otNavy] (1.5,2.35) circle (1.6pt);
  \fill[otNavy] (1.5,0.65) circle (1.6pt);
  \draw[otNavy,line width=0.8pt] (0,1.5) -- (1.5,2.35);
  \draw[otNavy,line width=0.8pt] (0,1.5) -- (1.5,0.65);
  \foreach \y/\lab in {2.35/,0.65/}{%
    \draw[otGreenDark,line width=0.7pt] (1.5,\y) -- (3,{\y+0.45});
    \draw[otGreenDark,line width=0.7pt] (1.5,\y) -- (3,{\y-0.45});
    \fill[otGreenDark] (3,{\y+0.45}) circle (1.3pt);
    \fill[otGreenDark] (3,{\y-0.45}) circle (1.3pt);}
  \node[font=\tiny,text=otNavy,above,sloped] at (0.72,1.98) {$P(B)$};
  \node[font=\tiny,text=otNavy,below,sloped] at (0.72,1.02) {$P(\overline B)$};
  \node[font=\tiny,text=otGreenDark,above,sloped] at (2.3,2.64) {$P(A\mid B)$};
  \node[font=\tiny,text=otGreenDark,anchor=west] at (3.08,2.8) {$A$};
  \node[font=\tiny,text=otGreenDark,anchor=west] at (3.08,1.9) {$\overline A$};
  \node[font=\tiny,text=otGreenDark,anchor=west] at (3.08,1.1) {$A$};
  \node[font=\tiny,text=otGreenDark,anchor=west] at (3.08,0.2) {$\overline A$};
  \node[font=\tiny,text=otNavy,anchor=east] at (1.42,2.6) {$B$};
  \node[font=\tiny,text=otNavy,anchor=east] at (1.42,0.4) {$\overline B$};
\end{tikzpicture}\\[0.5mm]
\legende{L'arbre pondéré : on multiplie le long des branches
(composées), on somme les feuilles d'un même événement (totales).}
\end{minipage}

\essSection{4}{Lois usuelles ; couples de variables}

\begin{definition}[Loi d'une variable aléatoire]
La \textbf{loi} de $X$ est la probabilité $P_X$ sur $X(\Omega)$ déterminée
par la distribution $\bigl(P(X=x)\bigr)_{x}$ ; on note $X\sim Y$ l'égalité
des lois, et alors $f(X)\sim f(Y)$. Lois à connaître :
\begin{itemize}[label={\color{otNavy}\textbullet},leftmargin=1.4em,itemsep=0.6mm,topsep=0.8mm]
  \item \textbf{uniforme} $\mathcal{U}(E)$ :
    $P(X=x)=\dfrac{1}{|E|}$ pour tout $x\in E$ ;
  \item \textbf{Bernoulli} $\mathcal{B}(p)$ :
    $P(X=1)=p$, $P(X=0)=1-p$ — un succès ou un échec ; l'indicatrice
    $\ind{A}$ suit $\mathcal{B}\bigl(P(A)\bigr)$ ;
  \item \textbf{binomiale} $\mathcal{B}(n,p)$ :
    $P(X=k)=\dbinom nk p^{k}(1-p)^{n-k}$ — le nombre de succès de $n$
    répétitions indépendantes.
\end{itemize}
Le \textbf{couple} $(X,Y)$ est une variable à valeurs dans un produit ; sa
loi \emph{conjointe} $P(X=x,Y=y)$ détermine les lois \emph{marginales}
(sommer sur l'autre variable), mais la réciproque est fausse
(exercice 10). On définit de même la loi conditionnelle de $X$ sachant un
événement, et les $n$-uplets.
\end{definition}

\essSection{5}{Indépendance}

\begin{definition}[Événements et variables indépendants]
$A$ et $B$ sont \textbf{indépendants} si $P(A\cap B)=P(A)P(B)$ — soit
$P(A\mid B)=P(A)$ quand $P(B)>0$ : savoir $B$ n'apprend rien sur $A$. Une
famille finie est indépendante si la formule produit vaut pour
\emph{toute} sous-famille. Si $A$ et $B$ sont indépendants, $A$ et
$\overline B$ aussi. Les variables $X$ et $Y$ sont indépendantes
(notation $X\perp\!\!\!\perp Y$) si
$P(X\in A,\ Y\in B)=P(X\in A)\,P(Y\in B)$ pour tous $A,B$, ce qui revient
à $P(X=x,Y=y)=P(X=x)P(Y=y)$ pour tous $x,y$.
\end{definition}

\begin{theoreme}[Les outils de l'indépendance]
Si $X_1,\dots,X_n$ sont indépendantes de loi $\mathcal{B}(p)$, alors
$X_1+\dots+X_n\sim\mathcal{B}(n,p)$ : la binomiale \emph{est} une somme de
Bernoulli. Si $X\perp\!\!\!\perp Y$, alors $f(X)\perp\!\!\!\perp g(Y)$.
Plus généralement (\emph{lemme des coalitions}) : si $X_1,\dots,X_n$ sont
indépendantes, alors $f(X_1,\dots,X_m)$ et $g(X_{m+1},\dots,X_n)$ le sont
aussi — et de même pour plus de deux coalitions.
\end{theoreme}

\begin{attention}
L'indépendance \emph{deux à deux} n'implique pas l'indépendance de la
famille : pour deux pièces équilibrées, « la première donne pile », « la
seconde donne pile » et « les deux donnent le même côté » sont deux à deux
indépendants, mais pas mutuellement (exercice 9). Vérifier la formule
produit pour \emph{toutes} les sous-familles.
\end{attention}

\essSection{6}{Espérance}

\begin{theoreme}[Espérance et ses propriétés]
L'\textbf{espérance} d'une variable réelle (ou complexe) est
\[
\mathrm{E}(X)=\sum_{x\in X(\Omega)}x\,P(X=x)
=\sum_{\omega\in\Omega}X(\omega)\,P(\ens{\omega}) :
\]
la moyenne des valeurs, pondérée par leurs probabilités — un indicateur de
\emph{position}. Elle est linéaire, positive, croissante, avec
$|\mathrm{E}(X)|\leqslant\mathrm{E}(|X|)$ ; $\mathrm{E}(c)=c$,
$\mathrm{E}(\ind A)=P(A)$,
$\mathrm{E}\bigl(\mathcal{B}(p)\bigr)=p$ et
$\mathrm{E}\bigl(\mathcal{B}(n,p)\bigr)=np$ (linéarité !). La variable
$X-\mathrm{E}(X)$ est dite \emph{centrée}.
\end{theoreme}

\begin{theoreme}[Transfert ; produit de variables indépendantes]
Pour toute fonction $f$ :
\[
\mathrm{E}\bigl(f(X)\bigr)=\sum_{x\in X(\Omega)}f(x)\,P(X=x)
\]
— inutile de connaître la loi de $f(X)$ ; la formule s'applique aussi aux
couples et aux $n$-uplets. Et si $X\perp\!\!\!\perp Y$ :
$\mathrm{E}(XY)=\mathrm{E}(X)\,\mathrm{E}(Y)$ (extension à $n$ variables
indépendantes) — la réciproque est fausse (exercice 14).
\end{theoreme}

\essSection{7}{Variance, écart type, covariance}

\begin{definition}[Dispersion]
La \textbf{variance} d'une variable réelle est
$\mathrm{V}(X)=\mathrm{E}\bigl((X-\mathrm{E}(X))^{2}\bigr)\geqslant0$, et
l'\textbf{écart type} $\sigma(X)=\sqrt{\mathrm{V}(X)}$ : des indicateurs
de \emph{dispersion}. Formules :
\[
\mathrm{V}(X)=\mathrm{E}\bigl(X^{2}\bigr)-\mathrm{E}(X)^{2}
\quad\text{(König-Huygens)},\qquad
\mathrm{V}(aX+b)=a^{2}\,\mathrm{V}(X) .
\]
Pour $\mathcal{B}(p)$ : $\mathrm{V}=p(1-p)$ ; pour $\mathcal{B}(n,p)$ :
$\mathrm{V}=np(1-p)$. Si $\sigma(X)>0$, la variable
$\dfrac{X-\mathrm{E}(X)}{\sigma(X)}$ est \emph{centrée réduite}.
\end{definition}

\begin{theoreme}[Covariance ; variance d'une somme]
$\operatorname{Cov}(X,Y)
=\mathrm{E}\bigl((X-\mathrm{E}X)(Y-\mathrm{E}Y)\bigr)
=\mathrm{E}(XY)-\mathrm{E}(X)\mathrm{E}(Y)$ ; elle est nulle pour des
variables indépendantes (on dit alors \emph{décorrélées} — la réciproque
est fausse). Et
\[
\mathrm{V}(X+Y)=\mathrm{V}(X)+\mathrm{V}(Y)+2\operatorname{Cov}(X,Y) :
\]
pour des variables deux à deux décorrélées, les variances s'ajoutent — on
retrouve ainsi $\mathrm{V}\bigl(\mathcal{B}(n,p)\bigr)=np(1-p)$.
\end{theoreme}

\essSection{8}{Inégalités et loi faible des grands nombres}

\begin{theoreme}[Markov et Bienaymé-Tchebychev]
Pour $X\geqslant0$ et $a>0$ :
$P(X\geqslant a)\leqslant\dfrac{\mathrm{E}(X)}{a}$ (Markov — appliquer la
croissance de l'espérance à $a\,\ind{(X\geqslant a)}\leqslant X$). En
l'appliquant à $(X-\mathrm{E}X)^{2}$ :
\[
P\bigl(|X-\mathrm{E}(X)|\geqslant\varepsilon\bigr)
\leqslant\frac{\mathrm{V}(X)}{\varepsilon^{2}}
\qquad\text{(Bienaymé-Tchebychev)} :
\]
une variable s'écarte rarement de son espérance de plus de quelques écarts
types.
\end{theoreme}

\begin{theoreme}[Loi faible des grands nombres]
Soient $X_1,\dots,X_n$ indépendantes de même loi, d'espérance $\mu$ et de
variance $\sigma^{2}$, et $M_n=\dfrac{X_1+\dots+X_n}{n}$ leur moyenne.
Alors $\mathrm{E}(M_n)=\mu$, $\mathrm{V}(M_n)=\dfrac{\sigma^{2}}{n}$ et,
pour tout $\varepsilon>0$ :
\[
P\bigl(|M_n-\mu|\geqslant\varepsilon\bigr)
\leqslant\frac{\sigma^{2}}{n\,\varepsilon^{2}}
\xrightarrow[n\to+\infty]{}0 :
\]
la moyenne empirique se \emph{concentre} autour de l'espérance. C'est la
justification de l'approche fréquentiste : les fréquences observées
convergent (en probabilité) vers les probabilités théoriques.
\end{theoreme}

%% =====================================================================
%%  EXERCICES
%% =====================================================================
\exercices
\rubrique{Applications}

\begin{exercice}[1]{Modéliser une expérience}
\begin{questions}
  \item On lance deux dés équilibrés. Décrire l'univers et calculer la
    probabilité que la somme vaille $7$.
  \item On lance trois fois une pièce équilibrée. Quelle est la
    probabilité d'obtenir exactement deux piles ?
\end{questions}
\end{exercice}

\begin{exercice}[1]{Les règles de calcul}
On donne $P(A)=\dfrac12$, $P(B)=\dfrac25$ et $P(A\cap B)=\dfrac15$.
Calculer $P(A\cup B)$, $P(A\setminus B)$ et
$P\bigl(\overline A\cap\overline B\bigr)$.
\end{exercice}

\begin{exercice}[1]{Tirages sans remise}
Une urne contient $3$ boules rouges et $2$ blanches ; on tire deux boules
successivement, sans remise.
\begin{questions}
  \item Calculer $P(\text{les deux sont rouges})$ par la formule des
    probabilités composées.
  \item Que vaut la probabilité que la seconde soit rouge sachant que la
    première l'est ?
\end{questions}
\end{exercice}

\begin{exercice}[1]{La formule des probabilités totales}
L'urne $A$ contient $2$ rouges et $1$ blanche, l'urne $B$ contient $1$
rouge et $2$ blanches. On choisit une urne au hasard, puis une boule.
Quelle est la probabilité de tirer une rouge ? Dessiner l'arbre.
\end{exercice}

\begin{exercice}[1]{La formule de Bayes}
Une maladie touche $1\,\%$ d'une population. Un test la détecte chez
$95\,\%$ des malades, mais se déclenche aussi chez $10\,\%$ des personnes
saines.
\begin{questions}
  \item Calculer la probabilité qu'un test soit positif.
  \item Un test est positif : quelle est la probabilité que la personne
    soit réellement malade ? Commenter ce résultat surprenant.
\end{questions}
\end{exercice}

\begin{exercice}[1]{Uniforme et Bernoulli}
\begin{questions}
  \item Soit $X$ le résultat d'un dé équilibré. Calculer
    $\mathrm{E}(X)$ et $\mathrm{V}(X)$.
  \item Soit $Y\sim\mathcal{B}(p)$. Retrouver $\mathrm{E}(Y)=p$ et
    $\mathrm{V}(Y)=p(1-p)$ ; pour quelle valeur de $p$ la variance
    est-elle maximale ?
\end{questions}
\end{exercice}

\begin{exercice}[1]{La loi binomiale}
Soit $X\sim\mathcal{B}\bigl(5,\frac13\bigr)$.
\begin{questions}
  \item Calculer $P(X=2)$, et vérifier que les probabilités
    $P(X=k)$ se somment à $1$.
  \item Donner $\mathrm{E}(X)$ et $\mathrm{V}(X)$.
  \item De quelle expérience $X$ est-il le modèle naturel ?
\end{questions}
\end{exercice}

\begin{exercice}[1]{Espérance : transfert et linéarité}
$X$ est le résultat d'un dé équilibré.
\begin{questions}
  \item Calculer $\mathrm{E}\bigl(X^{2}\bigr)$ par la formule de
    transfert. A-t-on $\mathrm{E}(X^{2})=\mathrm{E}(X)^{2}$ ?
  \item Calculer $\mathrm{E}(2X+3)$.
  \item Justifier $\mathrm{E}(\ind A)=P(A)$ : l'espérance généralise la
    probabilité.
\end{questions}
\end{exercice}

\rubrique{Entraînement}

\begin{exercice}[2]{Indépendance : exemples et pièges}
\begin{questions}
  \item On lance deux dés. Montrer que « le premier dé est pair » et « la
    somme est paire » sont indépendants.
  \item On lance deux pièces. Soient $A$ : « la première donne pile »,
    $B$ : « la seconde donne pile », $C$ : « les deux donnent le même
    côté ». Montrer que $A,B,C$ sont deux à deux indépendants, mais pas
    mutuellement indépendants.
\end{questions}
\end{exercice}

\begin{exercice}[2]{Loi conjointe et lois marginales}
\begin{questions}
  \item Soit $(X,Y)$ de loi conjointe uniforme sur
    $\ens{0,1}^{2}$ ($\frac14$ partout). Déterminer les marginales ;
    $X$ et $Y$ sont-elles indépendantes ?
  \item Soit maintenant $P(X=Y=0)=P(X=Y=1)=\frac12$. Vérifier que les
    marginales sont les mêmes qu'en a), mais que $X$ et $Y$ ne sont plus
    indépendantes : la loi conjointe contient strictement plus
    d'information.
\end{questions}
\end{exercice}

\begin{exercice}[2]{Minimum et maximum de deux dés}
On lance deux dés indépendants $X,Y$ et on pose $U=\min(X,Y)$,
$V=\max(X,Y)$.
\begin{questions}
  \item Calculer $P(U\geqslant k)$ pour $1\leqslant k\leqslant6$, et en
    déduire la loi de $U$.
  \item Montrer que $\mathrm{E}(U)=\dfrac{91}{36}$, puis, sans calculer la
    loi de $V$, que $\mathrm{E}(V)=\dfrac{161}{36}$.
\end{questions}
\end{exercice}

\begin{exercice}[2]{Variance d'une somme, covariance}
Avec les notations de l'exercice 11 :
\begin{questions}
  \item Vérifier $\mathrm{E}(XY)=\mathrm{E}(X)\mathrm{E}(Y)$ et calculer
    $\mathrm{V}(X+Y)$.
  \item Montrer que $UV=XY$, et en déduire
    $\operatorname{Cov}(U,V)=\mathrm{E}(X)\mathrm{E}(Y)-\mathrm{E}(U)\mathrm{E}(V)$ ;
    vérifier qu'elle est strictement positive et interpréter.
\end{questions}
\end{exercice}

\begin{exercice}[2]{La binomiale comme somme de Bernoulli}
Soient $X_1,\dots,X_n$ indépendantes de loi $\mathcal{B}(p)$, et
$S=X_1+\dots+X_n$.
\begin{questions}
  \item Justifier que $S\sim\mathcal{B}(n,p)$ en comptant les issues.
  \item En déduire $\mathrm{E}(S)=np$ et $\mathrm{V}(S)=np(1-p)$ sans
    aucun calcul de somme.
  \item Si $X\sim\mathcal{B}(n,p)$ et $Y\sim\mathcal{B}(m,p)$ sont
    indépendantes, quelle est la loi de $X+Y$ ?
\end{questions}
\end{exercice}

\begin{exercice}[2]{Décorrélées mais pas indépendantes}
Soit $X$ uniforme sur $\ens{-1,0,1}$ et $Y=X^{2}$.
\begin{questions}
  \item Calculer $\operatorname{Cov}(X,Y)$.
  \item Montrer que $X$ et $Y$ ne sont pas indépendantes.
  \item Montrer en revanche que pour des \emph{indicatrices},
    $\operatorname{Cov}(\ind A,\ind B)=P(A\cap B)-P(A)P(B)$ : la
    covariance nulle équivaut alors à l'indépendance.
\end{questions}
\end{exercice}

\begin{exercice}[2]{Points fixes d'une permutation aléatoire}
Soit $\sigma$ uniforme sur $S_n$ et $X$ son nombre de points fixes
(chapitres 14 et 16).
\begin{questions}
  \item Écrire $X=\sum_{i=1}^{n}\ind{(\sigma(i)=i)}$ et en déduire
    $\mathrm{E}(X)=1$.
  \item Calculer $\mathrm{E}\bigl(X(X-1)\bigr)$ en comptant les couples de
    points fixes, et en déduire $\mathrm{V}(X)=1$.
\end{questions}
\end{exercice}

\begin{exercice}[2]{Markov et Tchebychev : premiers usages}
\begin{questions}
  \item Un serveur traite en moyenne $2$ requêtes par seconde
    ($X\geqslant0$). Majorer $P(X\geqslant10)$.
  \item Pour un dé équilibré, majorer
    $P\bigl(|X-3{,}5|\geqslant2\bigr)$ par Tchebychev, puis calculer la
    probabilité exacte. Que dire de la précision de la borne ?
\end{questions}
\end{exercice}

\begin{exercice}[2]{La loi faible des grands nombres}
On lance $n$ fois une pièce équilibrée et on note $M_n$ la fréquence des
piles.
\begin{questions}
  \item Donner $\mathrm{E}(M_n)$ et $\mathrm{V}(M_n)$.
  \item Montrer que
    $P\bigl(|M_n-\tfrac12|\geqslant\varepsilon\bigr)
    \leqslant\dfrac{1}{4n\varepsilon^{2}}$.
  \item Quelle valeur de $n$ garantit
    $P\bigl(|M_n-\tfrac12|\geqslant0{,}1\bigr)\leqslant5\,\%$ ?
\end{questions}
\end{exercice}

\begin{exercice}[2]{Le paradoxe des anniversaires}
On suppose les $365$ jours de naissance équiprobables et indépendants.
\begin{questions}
  \item Montrer que la probabilité que $n$ personnes aient toutes des
    anniversaires distincts vaut
    $\displaystyle\prod_{k=0}^{n-1}\frac{365-k}{365}$.
  \item En déduire qu'à partir de $n=23$ personnes, une coïncidence est
    plus probable qu'improbable. Commenter.
\end{questions}
\end{exercice}

\rubrique{Approfondissement}

\begin{exercice}[3]{Dérangements : la loi des points fixes}
Soit $\sigma$ uniforme sur $S_n$ et $X$ son nombre de points fixes.
\begin{questions}
  \item Montrer que $P(X=k)=\dfrac{1}{k!}\cdot\dfrac{D_{n-k}}{(n-k)!}$,
    où $D_m$ compte les dérangements (chapitre 16).
  \item En déduire que, à $k$ fixé,
    $P(X=k)\xrightarrow[n\to+\infty]{}\dfrac{\mathrm{e}^{-1}}{k!}$ :
    la loi limite des points fixes (une loi de Poisson avant l'heure),
    cohérente avec $\mathrm{E}(X)=\mathrm{V}(X)=1$.
\end{questions}
\end{exercice}

\begin{exercice}[3]{Concentration de la binomiale et sondages}
Soit $S\sim\mathcal{B}(n,p)$ et $F=\dfrac Sn$ la fréquence.
\begin{questions}
  \item Montrer que $p(1-p)\leqslant\dfrac14$, puis que
    $P\bigl(|F-p|\geqslant\varepsilon\bigr)
    \leqslant\dfrac{1}{4n\varepsilon^{2}}$ — quelle que soit la valeur
    (inconnue !) de $p$.
  \item Un sondage vise une précision de $5$ points avec un risque d'au
    plus $5\,\%$ : quelle taille d'échantillon la borne garantit-elle ?
\end{questions}
\end{exercice}

\begin{exercice}[3]{Le problème de Monty Hall}
Derrière une porte sur trois se cache une voiture. Le candidat choisit une
porte ; l'animateur, qui sait tout, ouvre une \emph{autre} porte, vide,
puis propose de changer.
\begin{questions}
  \item Modéliser soigneusement (position de la voiture, choix, porte
    ouverte) et calculer la probabilité de gagner en gardant sa porte,
    puis en changeant.
  \item Expliquer, avec un conditionnement, pourquoi l'intuition
    « 50-50 » est fausse.
\end{questions}
\end{exercice}

\begin{exercice}[3]{L'espérance du maximum}
Soient $X_1,\dots,X_n$ indépendantes, uniformes sur $\ens{1,\dots,N}$, et
$M=\max(X_1,\dots,X_n)$.
\begin{questions}
  \item Montrer que $P(M\leqslant k)=\bigl(\tfrac kN\bigr)^{n}$, puis que
    $\mathrm{E}(M)=\displaystyle\sum_{k=0}^{N-1}
    \Bigl(1-\bigl(\tfrac kN\bigr)^{n}\Bigr)$.
  \item Vérifier pour $n=2$, $N=6$ la cohérence avec l'exercice 11.
  \item À l'aide d'une somme de Riemann (chapitre 15), montrer que
    $\dfrac{\mathrm{E}(M)}{N}\xrightarrow[N\to+\infty]{}\dfrac{n}{n+1}$.
\end{questions}
\end{exercice}

\begin{exercice}[3]{La marche aléatoire symétrique}
Soient $\varepsilon_1,\dots,\varepsilon_n$ indépendantes, uniformes sur
$\ens{-1,1}$, et $S_n=\varepsilon_1+\dots+\varepsilon_n$.
\begin{questions}
  \item Calculer $\mathrm{E}(S_n)$ et $\mathrm{V}(S_n)$ ; que dit
    Tchebychev sur l'ordre de grandeur de $|S_n|$ ?
  \item Montrer que $P(S_{2n}=0)=\dbinom{2n}{n}4^{-n}$, et donner son
    équivalent $\dfrac{1}{\sqrt{\pi n}}$ (chapitres 11, 15 et 16) :
    le retour en $0$ devient rare, mais lentement.
\end{questions}
\end{exercice}

\begin{exercice}[3]{Weierstrass par Bernstein}
Pour $f$ continue sur $\intFF{0}{1}$ et $x\in\intFF{0}{1}$, on pose
$B_n(f)(x)=\mathrm{E}\bigl(f(S_n/n)\bigr)
=\displaystyle\sum_{k=0}^{n}f\bigl(\tfrac kn\bigr)\binom nk
x^{k}(1-x)^{n-k}$, où $S_n\sim\mathcal{B}(n,x)$.
\begin{questions}
  \item Vérifier que $B_n$ transforme les fonctions $1$ et
    $t\mapsto t$ en elles-mêmes.
  \item Soit $\varepsilon>0$ et $\delta$ un module de continuité uniforme
    de $f$ (chapitre 15). En séparant les $k$ selon que
    $\bigl|\tfrac kn-x\bigr|<\delta$ ou non, et en appliquant
    Bienaymé-Tchebychev à $S_n$, montrer que
    \[
    \bigl|B_n(f)(x)-f(x)\bigr|
    \leqslant\varepsilon+\frac{2\max|f|}{4n\delta^{2}} .
    \]
  \item Conclure : toute fonction continue sur un segment est limite
    uniforme de polynômes (théorème de Weierstrass) — les probabilités
    démontrent un grand théorème d'analyse.
\end{questions}
\end{exercice}

% !TEX root = ../../main.tex
%% Défi du chapitre 17 — Le collectionneur de vignettes
%% Résultats vérifiés par verifications/defis_16_20.py.

\rubrique{Défi}

\begin{defi}{Le collectionneur de vignettes}{2 h}
Un album compte $n\geqslant2$ vignettes différentes. Chaque achat fournit
une vignette, uniformément au hasard et indépendamment des autres. Après
$m$ achats, on modélise l'expérience par $\Omega=\intFF1n^{m}$ muni de la
probabilité uniforme. Pour $i\in\intFF1n$, on note $A_i$ l'événement « la
vignette $i$ n'a pas été obtenue », $Y_m=\sum_{i=1}^{n}\ind{A_i}$ le nombre
de vignettes manquantes, et $X_m=n-Y_m$ le nombre de vignettes distinctes
obtenues. Combien d'achats faut-il pour compléter l'album ?

\partieExo{A}{Le nombre moyen de vignettes obtenues}
\begin{questions}
  \item Calculer $P(A_i)$, puis montrer que
    $\mathrm{E}(X_m)=n\Bigl(1-\bigl(1-\frac1n\bigr)^{m}\Bigr)$.
  \item Montrer que $\mathrm{E}(Y_m)\leqslant n\,\mathrm{e}^{-m/n}$. Pour
    $n=100$, combien d'achats faut-il pour obtenir \emph{en moyenne} $90$
    vignettes différentes ?
\end{questions}

\partieExo{B}{Le seuil $n\ln n$}
\begin{questions}[resume]
  \item Montrer que $P(\text{album incomplet})=P(Y_m\geqslant1)\leqslant\mathrm{E}(Y_m)$.
    En déduire que si $m\geqslant n\ln n+cn$ avec $c>0$, l'album est
    complet avec une probabilité au moins $1-\mathrm{e}^{-c}$. Application
    numérique : $n=100$, $c=3$.
  \item Calculer $P(A_i\cap A_j)$ pour $i\neq j$, puis $\mathrm{V}(Y_m)$ ;
    montrer que $\mathrm{V}(Y_m)\leqslant\mathrm{E}(Y_m)$.
  \item À l'aide de Bienaymé-Tchebychev, montrer que
    $P(Y_m=0)\leqslant\dfrac{1}{\mathrm{E}(Y_m)}$. En utilisant
    $\ln(1-x)\geqslant-\frac{x}{1-x}$ sur $\intFO01$, en déduire que si
    $m\leqslant(n-1)(\ln n-c)$, l'album est complet avec une probabilité au
    plus $\mathrm{e}^{-c}$. Conclure : on assiste à un \emph{seuil} brutal
    autour de $n\ln n$ achats.
\end{questions}

\medskip
\begin{center}
\begin{tikzpicture}[x=0.058cm,y=3.2cm,line cap=round]
  \draw[otRule,line width=0.3pt] (0,0.5) -- (162,0.5) (0,1) -- (162,1);
  \draw[otMuted,line width=0.5pt,-{Stealth[length=1.6mm]}] (0,0) -- (168,0);
  \draw[otMuted,line width=0.5pt,-{Stealth[length=1.6mm]}] (0,0) -- (0,1.12);
  \draw[otGreenDark,line width=0.7pt,dash pattern=on 2.5pt off 1.5pt] (59.9,0) -- (59.9,1.05);
  \draw[otNavy!60,line width=0.7pt,dash pattern=on 2.5pt off 1.5pt] (119.9,0) -- (119.9,1.05);
  \draw[otNavy,line width=1.3pt] (20,0.0000) -- (25,0.0000) -- (30,0.0013) -- (35,0.0098) --
    (40,0.0359) -- (45,0.0875) -- (50,0.1642) -- (55,0.2585) -- (60,0.3606) -- (65,0.4615) --
    (70,0.5551) -- (75,0.6379) -- (80,0.7087) -- (85,0.7678) -- (90,0.8162) -- (95,0.8553) --
    (100,0.8865) -- (105,0.9113) -- (110,0.9308) -- (115,0.9462) -- (120,0.9582) --
    (125,0.9675) -- (130,0.9748) -- (135,0.9805) -- (140,0.9849) -- (145,0.9883) --
    (150,0.9909) -- (155,0.9930) -- (160,0.9946);
  \foreach \v in {0,40,80,120,160}{\node[font=\tiny,text=otMuted,below] at (\v,0) {$\v$};}
  \node[font=\tiny,text=otMuted,anchor=east] at (0,1) {$1$};
  \node[font=\tiny,text=otMuted,anchor=east] at (0,0.5) {$\tfrac12$};
  \node[font=\scriptsize,text=otGreenDark,anchor=south] at (59.9,1.05) {$n\ln n$};
  \node[font=\scriptsize,text=otNavy!70,anchor=south] at (119.9,1.05) {$n\ln n+3n$};
  \node[font=\scriptsize,text=otMuted,anchor=west] at (168,0) {$m$};
\end{tikzpicture}\\[0.5mm]
\legende{Pour $n=20$ vignettes : probabilité exacte d'avoir complété
l'album après $m$ achats (question 6). Elle passe de presque $0$ à presque
$1$ autour de $n\ln n\approx60$.}
\end{center}

\partieExo{C}{Formule exacte pour un petit album}
\begin{questions}[resume]
  \item À l'aide du défi du chapitre 16, montrer que
    $P(\text{album complet après }m\text{ achats})=\dfrac{n!\,S(m,n)}{n^{m}}$.
  \item Le cas du dé ($n=6$) : sachant que $S(12,6)=1\,323\,652$ et
    $S(13,6)=9\,321\,312$, déterminer le nombre minimal de lancers pour que
    la probabilité d'avoir obtenu les six faces dépasse $\frac12$. Comparer
    avec $n\ln n$.
\end{questions}
\end{defi}

